% !TEX root =./ECE343_Practice_main.tex
%
% L28 practice set (Magnetic Fields and Gauss's Law for Magnetism): a finite
% wire, a loop on its axis, a square loop, which fields can be magnetic, and
% the force on a moving charge.

\fullwidth{\textbf{Learning Objectives 7.1 and 7.2: } I can apply the Biot-Savart law to
simple magnetostatic systems to find the magnetic field at an observation point.
I can articulate the meaning of both the integral form and point form of
Gauss's Law for magnetism.
\\[4pt]\normalfont\small Show your work. Take $\mu = \mu_0 = 4\pi\ee{-7}$ H/m. Write a
\textbf{Documentation} line at the end (\textbf{None} if you did not collaborate).}

\begin{questions}

%%%%%%%%%%%%%%%%%%%%%%%%%%%%%%%%%%%%%%%%%
\question \textbf{A finite wire.} A straight wire on the $z$ axis carries $I = 2$ A in
the $+z$ direction from $z = -10$ cm to $z = 10$ cm.
\begin{parts}

	\part Starting from the Biot-Savart law, derive the field at a distance $\rho$
	from the wire in the plane $z = 0$.
	\begin{solution}[2.6in]
	With $d\mathbf{l} = \hat{\mathbf{z}}\ dz^{\prime}$ and
	$\mathbf{R} = \rho\hat{\boldsymbol{\rho}} - z^{\prime}\hat{\mathbf{z}}$,
	the cross product is
	$d\mathbf{l}\times\mathbf{R} = \hat{\boldsymbol{\phi}}\ \rho\ dz^{\prime}$, so
	\eq{
		\mathbf{H} &= \int_{-L/2}^{L/2} \dfrac{I\ \hat{\boldsymbol{\phi}}\ \rho\ dz^{\prime}}{4\pi\lp\rho^2 + z^{\prime 2}\rp^{3/2}}
		= \hat{\boldsymbol{\phi}}\dfrac{I\rho}{4\pi}\lb \dfrac{z^{\prime}}{\rho^2\sqrt{\rho^2 + z^{\prime 2}}} \rb_{-L/2}^{L/2} \\
		&= \hat{\boldsymbol{\phi}}\dfrac{I L}{4\pi\rho\sqrt{\rho^2 + L^2/4}}
	}
	The integral follows from $z^{\prime} = \rho\tan\alpha$, as in Lesson 18.
	\end{solution}

	\begin{minipage}[t]{0.58\linewidth}
		\part Find $\mathbf{H}$ and $|\mathbf{B}|$ at the point $(5, 0, 0)$ cm.
		\begin{solution}[1.3in]
		\eq{
			H &= \dfrac{(2)(0.2)}{4\pi(0.05)\sqrt{0.05^2 + 0.1^2}} = 5.69\ \text{A/m} \\
			\mathbf{H} &= 5.69\ \hat{\mathbf{y}}\ \text{A/m} \\
			B &= \mu_0 H = 7.16\ \mu\text{T}
		}
		At $(5, 0, 0)$ cm, $\hat{\boldsymbol{\phi}} = \hat{\mathbf{y}}$.
		\end{solution}
	\end{minipage}\hfill%
	\begin{minipage}[t]{0.37\linewidth}
		\raggedleft \vspace{12pt}
		$\mathbf{H}$ = \ansbox[1.5in]{$5.69\ \hat{\mathbf{y}}$ A/m}

		\vspace{4pt}
		$|\mathbf{B}|$ = \ansbox[1.5in]{$7.16\ \mu$T}
	\end{minipage}

	\begin{minipage}[t]{0.58\linewidth}
		\part Compare your answer with the field of an infinite wire carrying the same
		current. What fraction of the infinite-wire field does the finite wire give?
		\begin{solution}[1in]
		\eq{
			H_\infty &= \dfrac{I}{2\pi\rho} = \dfrac{2}{2\pi(0.05)} = 6.37\ \text{A/m} \\
			\dfrac{H}{H_\infty} &= \dfrac{5.69}{6.37} = 0.894
		}
		\end{solution}
	\end{minipage}\hfill%
	\begin{minipage}[t]{0.37\linewidth}
		\raggedleft \vspace{12pt}
		fraction = \ansbox[1.3in]{$0.894$}
	\end{minipage}

\end{parts}

\newpage
%%%%%%%%%%%%%%%%%%%%%%%%%%%%%%%%%%%%%%%%%
\question \textbf{A loop on its axis.} A circular loop of radius $a = 4$ cm lies in the
plane $z = 0$, centered on the $z$ axis, and carries $I = 0.5$ A counterclockwise
seen from above.
\begin{parts}

	\part Starting from the Biot-Savart law, derive the field at height $z$ on the
	axis. Explain why the radial part of the integral is zero.
	\begin{solution}[2.6in]
	With $d\mathbf{l} = \hat{\boldsymbol{\phi}}^{\prime}\ a\ d\phi^{\prime}$ and
	$\mathbf{R} = z\hat{\mathbf{z}} - a\hat{\boldsymbol{\rho}}^{\prime}$,
	$R = \sqrt{a^2 + z^2}$ for every element and
	$d\mathbf{l}\times\mathbf{R} = a\ d\phi^{\prime}\lp z\hat{\boldsymbol{\rho}}^{\prime} + a\hat{\mathbf{z}}\rp$:
	\eq{
		\mathbf{H} &= \dfrac{I a}{4\pi\lp a^2 + z^2\rp^{3/2}}\int_0^{2\pi}\lp z\hat{\boldsymbol{\rho}}^{\prime} + a\hat{\mathbf{z}}\rp d\phi^{\prime}
		= \hat{\mathbf{z}}\dfrac{I a^2}{2\lp a^2 + z^2\rp^{3/2}}
	}
	The unit vector $\hat{\boldsymbol{\rho}}^{\prime}$ turns with $\phi^{\prime}$, and its
	integral around a full circle is zero: the elements on opposite sides of the
	loop produce equal and opposite radial fields.
	\end{solution}

	\begin{minipage}[t]{0.58\linewidth}
		\part Find $|\mathbf{H}|$ at the center of the loop and at $z = 3$ cm.
		\begin{solution}[1.2in]
		\eq{
			H(0) &= \dfrac{I}{2a} = \dfrac{0.5}{2(0.04)} = 6.25\ \text{A/m} \\
			H(3\ \text{cm}) &= \dfrac{(0.5)(0.04)^2}{2\lp 0.04^2 + 0.03^2\rp^{3/2}} = 3.20\ \text{A/m}
		}
		\end{solution}
	\end{minipage}\hfill%
	\begin{minipage}[t]{0.37\linewidth}
		\raggedleft \vspace{12pt}
		center = \ansbox[1.3in]{$6.25$ A/m}

		\vspace{4pt}
		$z = 3$ cm: \ansbox[1.3in]{$3.20$ A/m}
	\end{minipage}

	\begin{minipage}[t]{0.58\linewidth}
		\part At what height on the axis is the field half its value at the center?
		\begin{solution}[1.2in]
		\eq{
			\lp\dfrac{a^2}{a^2 + z^2}\rp^{3/2} &= \dfrac{1}{2} \implies a^2 + z^2 = 2^{2/3}a^2 \\
			z &= a\sqrt{2^{2/3} - 1} = 0.766\ a = 3.07\ \text{cm}
		}
		\end{solution}
	\end{minipage}\hfill%
	\begin{minipage}[t]{0.37\linewidth}
		\raggedleft \vspace{12pt}
		$z$ = \ansbox[1.3in]{$3.07$ cm}
	\end{minipage}

\end{parts}

\newpage
%%%%%%%%%%%%%%%%%%%%%%%%%%%%%%%%%%%%%%%%%
\question \textbf{A square loop.} A square loop of side $s = 8$ cm carries $I = 2$ A.
\begin{parts}

	\begin{minipage}[t]{0.58\linewidth}
		\part Use the finite-wire result to find the field at the center of the loop.
		\begin{solution}[2in]
		The center is on the bisecting plane of each side, with $L = s$ and
		$\rho = s/2$, and all four sides produce fields in the same direction:
		\eq{
			H_\text{side} &= \dfrac{I s}{4\pi(s/2)\sqrt{s^2/4 + s^2/4}} = \dfrac{\sqrt{2}\ I}{2\pi s} = 5.63\ \text{A/m} \\
			H &= 4H_\text{side} = \dfrac{2\sqrt{2}\ I}{\pi s} = 22.5\ \text{A/m}
		}
		\end{solution}
	\end{minipage}\hfill%
	\begin{minipage}[t]{0.37\linewidth}
		\raggedleft \vspace{12pt}
		$H$ = \ansbox[1.3in]{$22.5$ A/m}
	\end{minipage}

	\part A circular loop of radius $a = s/2$ fits inside the square and carries the
	same current. Find its field at the center, and explain why it is larger than
	the square's even though the square has more wire.
	\begin{solution}[1.3in]
	\eq{ H = \dfrac{I}{2a} = \dfrac{2}{2(0.04)} = 25.0\ \text{A/m} }
	Every element of the circle is $4$ cm from the center and broadside to it. The
	wire near the corners of the square is farther away, up to
	$s/\sqrt{2} = 5.66$ cm, and sees the center at a smaller $\sin\theta$, so it
	contributes less.
	\end{solution}

\end{parts}

\newpage
%%%%%%%%%%%%%%%%%%%%%%%%%%%%%%%%%%%%%%%%%
\question \textbf{Gauss's law for magnetism.}
\begin{parts}

	\part Which of these fields could be a magnetic flux density? Justify each answer.
	\[ \mathbf{B}_1 = x\hat{\mathbf{x}} - y\hat{\mathbf{y}}, \qquad \mathbf{B}_2 = y\hat{\mathbf{x}} + x\hat{\mathbf{y}}, \qquad \mathbf{B}_3 = x\hat{\mathbf{x}} + y\hat{\mathbf{y}} \]
	\begin{solution}[1.3in]
	\eq{
		\nabla\cdot\mathbf{B}_1 &= 1 - 1 = 0 \quad \text{possible} \\
		\nabla\cdot\mathbf{B}_2 &= 0 + 0 = 0 \quad \text{possible} \\
		\nabla\cdot\mathbf{B}_3 &= 1 + 1 = 2 \quad \text{not possible: a magnetic field has } \nabla\cdot\mathbf{B} = 0
	}
	\end{solution}

	\begin{minipage}[t]{0.58\linewidth}
		\part The Earth's field at a site is $50\ \mu$T, pointing straight down. Find the
		magnetic flux through a hemispherical dome of radius $1$ m sitting on flat
		ground. Use Gauss's law to avoid integrating over the dome.
		\begin{solution}[1.6in]
		The dome and the flat disk under it form a closed surface, so the flux
		through the dome equals the flux through the disk:
		\eq{ \Phi = B\pi R^2 = \lp 50\ee{-6}\rp\pi\lp 1\rp^2 = 157\ \mu\text{Wb} }
		\end{solution}
	\end{minipage}\hfill%
	\begin{minipage}[t]{0.37\linewidth}
		\raggedleft \vspace{12pt}
		$\Phi$ = \ansbox[1.3in]{$157\ \mu$Wb}
	\end{minipage}

\end{parts}

\newpage
%%%%%%%%%%%%%%%%%%%%%%%%%%%%%%%%%%%%%%%%%
\question \textbf{The force on a moving charge.} A long straight wire on the $z$ axis
carries $10$ A in the $+z$ direction. An electron passes the point $(1, 0, 0)$ cm
moving at $10^6$ m/s in the $+z$ direction.
\begin{parts}

	\begin{minipage}[t]{0.58\linewidth}
		\part Find $\mathbf{B}$ at the electron.
		\begin{solution}[1in]
		\eq{ \mathbf{B} = \hat{\boldsymbol{\phi}}\dfrac{\mu_0 I}{2\pi\rho} = \hat{\mathbf{y}}\dfrac{\lp 4\pi\ee{-7}\rp(10)}{2\pi(0.01)} = 200\ \hat{\mathbf{y}}\ \mu\text{T} }
		\end{solution}
	\end{minipage}\hfill%
	\begin{minipage}[t]{0.37\linewidth}
		\raggedleft \vspace{12pt}
		$\mathbf{B}$ = \ansbox[1.3in]{$200\ \hat{\mathbf{y}}\ \mu$T}
	\end{minipage}

	\begin{minipage}[t]{0.58\linewidth}
		\part Find the force on the electron, and state whether it pulls the electron
		toward the wire or pushes it away.
		\begin{solution}[1.4in]
		\eq{
			\mathbf{F} &= q\mathbf{v}\times\mathbf{B} = \lp -1.602\ee{-19}\rp\lp 10^6\hat{\mathbf{z}}\rp\times\lp 2\ee{-4}\hat{\mathbf{y}}\rp \\
			&= \lp -3.20\ee{-17}\rp\lp -\hat{\mathbf{x}}\rp = 3.20\ee{-17}\ \hat{\mathbf{x}}\ \text{N}
		}
		The force points along $+\hat{\mathbf{x}}$, away from the wire: the cross
		product gives $-\hat{\mathbf{x}}$, and the electron's negative charge
		reverses it.
		\end{solution}
	\end{minipage}\hfill%
	\begin{minipage}[t]{0.37\linewidth}
		\raggedleft \vspace{12pt}
		$\mathbf{F}$ = \ansbox[1.5in]{$3.20\ee{-17}\ \hat{\mathbf{x}}$ N}
	\end{minipage}

\end{parts}

\vspace{1.5em}
\noindent\textbf{Documentation:} \rule[-2pt]{0.6\linewidth}{0.4pt}

\end{questions}
